% TOP_PROBLEM: {"id": 30007230, "title": "The Kerr–Newman area bound without electromagnetic-tail controls", "category_id": 16, "category": {"id": 16, "name": "physics", "display_name": "Mathematical Physics", "description": "Problems at the intersection of mathematics and physics.", "slug": "physics", "order_index": 16, "created_at": "2026-07-31T15:26:25.671Z"}, "status": "open"}
% FIRST_POSTED: 2026-10-06
% ATTEMPT_STATUS: solved
% OPENAI_CLAIM: {"schema_version":1,"source_repo":"https://github.com/openai/math","source_commit":"adc7f1241b42e322a6451854ab7e4b4c146bf78a","release_date":"2026-10-06","match":"full","resolution":"disproved","independently_reviewed":false,"second_pass":"agreed","families":[{"family":"260","title":"Spacetime Penrose inequalities: enclosing area, charge, rotation, and anti-de Sitter extensions","subject":"Mathematical physics","manuscripts":[{"title":"Electromagnetic tails and the Kerr–Newman Penrose inequality","pdf_path":"preprints/Electromagnetic-tails-and-the-Kerr-Newman-Penrose-inequality-October-5-2026/paper.pdf","readme_path":"preprints/Electromagnetic-tails-and-the-Kerr-Newman-Penrose-inequality-October-5-2026/README.md","date":"2026-10-05","theorem_ref":"preprints/Electromagnetic-tails-and-the-Kerr-Newman-Penrose-inequality-October-5-2026/build/sections/00-introduction.tex:145 (thm:main)"}],"lean":{"doc_path":"lean/docs/260.md","comparators":[],"declarations":[],"covers_main_theorem":false},"reasoning_trace":null}]}
\documentclass[11pt]{article}
\usepackage[margin=25mm]{geometry}
\usepackage{fontspec}
\setmainfont{Latin Modern Roman}
\usepackage{amsmath,amssymb,mathtools}
\usepackage{unicode-math}
\setmathfont{Latin Modern Math}
\usepackage{enumitem}
\usepackage{xurl,hyperref,bookmark}
\hypersetup{colorlinks=true,urlcolor=blue}
\setcounter{secnumdepth}{0}
\setlength{\parindent}{0pt}
\setlength{\parskip}{5pt}
\setlength{\emergencystretch}{3em}
\begin{document}
\section*{The Kerr–Newman area bound without electromagnetic-tail controls}

\subsection{OpenAI's claim}
OpenAI presents a counterexample in the selected manuscript. With angularly varying $O(r^{-2})$ electromagnetic tails and bare gravitational angular momentum, the authors claim Penrose violations and non-Kerr--Newman equality examples.

\subsection{Outline of the argument}
Explicit axisymmetric electrovacuum seeds are reduced to a scalar conformal equation. Tail choices make both total charges zero while retaining nonzero fields and altering the angular-momentum balance; the conformal solution preserves the required outer horizon and realizes the inequality and rigidity failures.

\subsection{Scope relative to this problem}
The main theorem constructs admissible zero-total-charge exteriors with m\({}^{2}\) below the requested gravitational-J expression, and equality examples not induced from Kerr–Newman. Both universal inequality and equality-rigidity assertions are refuted. Within these assumptions, this is a claimed negative answer to the stated question.

\subsection{Formal verification}
For family 260, the only selected family entry is the CKS end-replacement construction and Schwarzschild equality examples. A connected-boundary inequality is conditional on the flat Penrose input; the general Bondi or spacetime inequalities are not formalized there. No Comparator is assigned to the selected manuscript here; the family page selects other papers or supporting statements. This site has not executed the Lean code. The scope descriptions do not certify the remaining manuscript arguments.

\subsection{Status caveat}
OpenAI attributes these manuscripts to work by an internal model. Its release cautions that arguments without formalization can contain mistakes. This catalogue has not independently checked the proofs; the labels report the mathematical claims attributed to OpenAI.

\subsection{Sources}
\begin{itemize}
\item OpenAI, \emph{Electromagnetic tails and the Kerr–Newman Penrose inequality}, 2026-10-05; citation key \texttt{OAI:Electromagnetic-tails-and-the-Kerr-Newman-Penrose-inequality-October-5-2026}. \href{https://github.com/openai/math/blob/adc7f1241b42e322a6451854ab7e4b4c146bf78a/preprints/Electromagnetic-tails-and-the-Kerr-Newman-Penrose-inequality-October-5-2026/paper.pdf}{PDF (GitHub)}; \href{https://github.com/openai/math/raw/adc7f1241b42e322a6451854ab7e4b4c146bf78a/preprints/Electromagnetic-tails-and-the-Kerr-Newman-Penrose-inequality-October-5-2026/paper.pdf}{PDF (direct)}; \href{https://github.com/openai/math/blob/adc7f1241b42e322a6451854ab7e4b4c146bf78a/preprints/Electromagnetic-tails-and-the-Kerr-Newman-Penrose-inequality-October-5-2026/README.md}{Manuscript page and citation}; \href{https://github.com/openai/math/blob/main/preprints/Electromagnetic-tails-and-the-Kerr-Newman-Penrose-inequality-October-5-2026/paper.pdf}{Latest version}.
\item Family 260: \href{https://github.com/openai/math/blob/adc7f1241b42e322a6451854ab7e4b4c146bf78a/lean/docs/260.md}{Lean scope}.
\item \href{https://github.com/openai/math/blob/adc7f1241b42e322a6451854ab7e4b4c146bf78a/CONTENTS.md}{OpenAI catalogue entry}.
\end{itemize}
\end{document}
